% ============================================================
%  SECTION 3: KIỂM THỬ
% ============================================================
\chapter{Kiểm thử}

\section{Kiểm thử CI Pipeline (Yêu cầu 3)}

\subsection{Kịch bản test}

\begin{lstlisting}[style=bashstyle, caption=Test CI pipeline]
# Tao branch moi
git checkout -b dev_tax_service

# Chinh sua code trong service tax
echo "// test change" >> tax/src/main/java/com/yas/tax/TaxApplication.java

# Commit va push
git add .
git commit -m "test: trigger CI pipeline for tax service"
git push origin dev_tax_service
\end{lstlisting}

\subsection{Kết quả mong đợi}

\begin{enumerate}
  \item Workflow \texttt{CI Build and Push} tự động trigger
  \item Job \texttt{detect-changes} phát hiện thư mục \texttt{tax/} thay đổi
  \item Job \texttt{build-java} chạy: Maven build \textrightarrow{} Docker build \textrightarrow{} push image
  \item Image được push lên Docker Hub: \texttt{<DOCKER\_USER>/tax:<commit\_sha>}
  \item Có thể kiểm tra tại: \href{https://hub.docker.com}{hub.docker.com}
\end{enumerate}

\begin{successbox}[Kết quả kiểm thử CI]
\begin{lstlisting}[style=bashstyle]
# Kiem tra image tren Docker Hub
docker pull <DOCKER_USER>/tax:abc1234
# Status: Downloaded newer image

# Xac nhan tag correct
docker image ls | grep tax
# <DOCKER_USER>/tax   abc1234   sha256:...   2 minutes ago
\end{lstlisting}
\end{successbox}

% ============================================================
\section{Kiểm thử Developer Build (Yêu cầu 4)}

\subsection{Kịch bản test}

\begin{enumerate}
  \item Vào GitHub \textrightarrow{} Actions \textrightarrow{} \textbf{developer\_build} \textrightarrow{} \textbf{Run workflow}
  \item Nhập:
        \begin{itemize}
          \item \texttt{service\_name}: \texttt{tax}
          \item \texttt{branch\_name}: \texttt{dev\_tax\_service}
        \end{itemize}
  \item Click \textbf{Run workflow}
\end{enumerate}

\subsection{Kết quả}

\begin{lstlisting}[style=bashstyle, caption=Kiểm tra sau khi developer\_build chạy]
# Kiem tra namespace developer
kubectl get pods -n developer
# NAME                              READY   STATUS    RESTARTS   AGE
# cart-xxx                          2/2     Running   0          2m
# order-xxx                         2/2     Running   0          2m
# tax-xxx                           2/2     Running   0          2m  <- custom image
# ...

# Lay NodePort cua tax service
kubectl get svc tax -n developer
# NAME   TYPE       CLUSTER-IP    EXTERNAL-IP   PORT(S)          AGE
# tax    NodePort   10.43.x.x     <none>        80:3xxxx/TCP     2m

# Them vao /etc/hosts tren may developer
<WORKER_IP> developer.yas.local.com api.developer.yas.local.com

# Test truy cap
curl http://api.developer.yas.local.com/tax/actuator/health
# {"status":"UP"}
\end{lstlisting}

% ============================================================
\section{Kiểm thử Cleanup (Yêu cầu 5)}

\begin{lstlisting}[style=bashstyle, caption=Test cleanup workflow]
# Sau khi chay cleanup_developer_build workflow
kubectl get namespace developer
# Error from server (NotFound): namespaces "developer" not found
# -> Namespace da bi xoa thanh cong
\end{lstlisting}

% ============================================================
\section{Kiểm thử Dev/Staging Deploy (Yêu cầu 6)}

\subsection{Test Dev Environment}

\begin{lstlisting}[style=bashstyle, caption=Trigger dev deploy bằng push vào main]
git checkout main
git merge dev_tax_service
git push origin main
# -> Workflow npt-dev-deploy.yml tu dong chay
# -> ArgoCD ApplicationSet sync -> deploy vao namespace dev
\end{lstlisting}

\begin{lstlisting}[style=bashstyle, caption=Kiểm tra namespace dev]
kubectl get pods -n dev | grep -c Running
# 20+  <- tat ca service dang Running

curl http://storefront.dev.yas.local.com
# HTTP 200
\end{lstlisting}

\subsection{Test Staging Environment}

\begin{lstlisting}[style=bashstyle, caption=Trigger staging deploy bằng tag]
git tag v1.0.0
git push origin v1.0.0
# -> Workflow npt-staging-deploy.yml chay
# -> Build tat ca services voi tag :v1.0.0
# -> Deploy vao namespace staging

kubectl get pods -n staging
# All running with image tag v1.0.0
\end{lstlisting}

% ============================================================
\section{Kiểm thử Istio mTLS (Nâng cao)}

\subsection{Kiểm tra mTLS đang hoạt động}

\begin{lstlisting}[style=bashstyle, caption=Verify mTLS STRICT mode]
# Kiem tra PeerAuthentication
kubectl get peerauthentication -n yas
# NAME           MODE     AGE
# default-mtls   STRICT   1d

# Kiem tra pod co 2/2 containers (co sidecar)
kubectl get pods -n yas | grep "2/2"
# Tat ca pods phai co 2/2 Running

# Verify bang istioctl
./istio-1.24.3/bin/istioctl experimental authz check \
  $(kubectl get pod -n yas -l app.kubernetes.io/name=product \
    -o jsonpath='{.items[0].metadata.name}') -n yas
\end{lstlisting}

\subsection{Kiểm tra Authorization Policy}

\begin{lstlisting}[style=bashstyle, caption=Test Authorization Policy]
# Test 1: order duoc phep goi tax (co trong policy)
kubectl exec -n yas deploy/order -- \
  curl -v http://tax/tax/actuator/health
# -> HTTP 200 OK  (THANH CONG - duoc phep)

# Test 2: search khong duoc phep goi tax (khong co trong policy)
kubectl exec -n yas deploy/search -- \
  curl -v http://tax/tax/actuator/health
# -> RBAC: access denied
# -> HTTP 403 Forbidden  (THANH CONG - bi chan)
\end{lstlisting}

\subsection{Kiểm tra Retry Policy}

\begin{lstlisting}[style=bashstyle, caption=Test Retry Policy]
# Xem log Envoy de thay retry
kubectl logs deploy/order -n yas -c istio-proxy | grep "retry\|upstream"
# upstream_request_count: 3  <- thu lai 3 lan

# Kiem tra config retry tren VirtualService
kubectl get virtualservice tax-retry -n yas -o yaml
\end{lstlisting}

\subsection{Kiali Topology}

\begin{lstlisting}[style=bashstyle, caption=Truy cập Kiali và kiểm tra topology]
# Port-forward Kiali
kubectl port-forward svc/kiali -n istio-system 20001:20001 &

# Truy cap: http://localhost:20001
# Menu: Graph -> Namespace: yas
# Bat: Display -> Security (hien thi lock icon cho mTLS)
# Ket qua: Tat ca service-to-service connection hien lock icon xanh = mTLS active
\end{lstlisting}

\begin{infobox}[Kết quả Kiali sau khi hoàn thiện cấu hình]
\begin{itemize}
  \item Tất cả 21 DestinationRule hiển thị \textbf{[OK] xanh} (không còn wildcard host)
  \item Tất cả 18 Service hiển thị \textbf{[OK] xanh} (đã fix port name \texttt{http-metric})
  \item AuthorizationPolicy hiển thị đúng workload (đã sửa selector sang \texttt{app.kubernetes.io/name})
  \item Service mesh topology hiển thị đầy đủ traffic flow với mTLS lock icon
\end{itemize}
\end{infobox}

% ============================================================
\section{Kiểm thử End-to-End}

\begin{lstlisting}[style=bashstyle, caption=Test tổng thể các URL hệ thống]
# Storefront (ung dung chinh)
curl -o /dev/null -w "%{http_code}" http://storefront.yas.local.com
# -> 200

# Backoffice
curl -o /dev/null -w "%{http_code}" http://backoffice.yas.local.com
# -> 302 (redirect den Keycloak login)

# Keycloak (Identity)
curl -o /dev/null -w "%{http_code}" http://identity.yas.local.com
# -> 302

# pgAdmin
curl -o /dev/null -w "%{http_code}" http://pgadmin.yas.local.com
# -> 302

# AKHQ (Kafka UI)
curl -o /dev/null -w "%{http_code}" http://akhq.yas.local.com
# -> 307

# Grafana (Observability)
curl -o /dev/null -w "%{http_code}" http://grafana.yas.local.com
# -> 302
\end{lstlisting}

\begin{lstlisting}[style=bashstyle, caption=Kiểm tra tổng thể pods]
kubectl get pods -A --no-headers | awk '{print $4}' | sort | uniq -c | sort -rn
# 56+ Running
#   3 Completed   <- ingress admission jobs, binh thuong
# 0 Error / CrashLoopBackOff / Pending
\end{lstlisting}
